\subsection{Autorización con ámbito y protección de datos}
\subsubsection{Autorización con ámbito y seguridad a nivel de campo}
\label{anx:autorizacion-datos}

El modelo de autorización de los microservicios anteriores representaba el rol como una
cadena de texto en la tabla de usuarios, y el control de acceso se reducía a verificar que
esa cadena perteneciera a un conjunto permitido. Dos usuarios con el mismo rol sobre
colegios distintos superaban, así, exactamente el mismo control. La única defensa consistía
en que cada punto de entrada recordara comparar el colegio del token con el del recurso
solicitado, y la auditoría de los 152 puntos de entrada de esos servicios encontró 22 en los
que esa comparación faltaba o era parcial (Sección~\ref{sec:defectos}).

El diagnóstico requiere precisión. Uno de esos servicios, el de instituciones, ya definía
una relación de pertenencia con la tripla \emph{(usuario, institución, rol)}, introducida
para corregir una referencia directa insegura real. Expresar autoridad con ámbito no era,
por tanto, imposible: el defecto era de alcance y de completitud. El mecanismo cubría un
servicio de siete; admitía un único nivel de ámbito, la institución, cuando el recurso
protegido pertenece a una sede y un grado; permitía un solo rol por institución, y carecía
de vigencia temporal. Esta última carencia producía una contradicción ilustrativa: la
documentación del modelo justificaba conservar el vínculo para auditar quién tuvo acceso y
cuándo, pero la única forma de revocarlo era eliminarlo. El diseño generaliza esa solución,
ya validada contra un ataque real, con la jerarquía de ámbitos, la vigencia temporal y la
aplicación uniforme a todos los servicios.

Con el control de acceso basado en roles con ámbito, la pregunta que el sistema responde
deja de ser \emph{qué es este usuario} y pasa a ser \emph{qué puede hacer este usuario sobre
este recurso concreto}:

\begin{equation}
\mathrm{puede}(u, p, a) \iff \exists\, r \in \mathcal{R}_{\text{vig}}(u, a) : (r, p) \in \mathcal{P}
\end{equation}

\noindent donde $\mathcal{P}$ es el conjunto de pares de rol y permiso, y
$\mathcal{R}_{\text{vig}}(u, a)$ el de roles vigentes del usuario $u$ sobre el ámbito $a$ o
sobre cualquier ámbito que lo contenga, según la jerarquía
$\text{global} \succ \text{colegio} \succ \text{sede} \succ \text{matrícula}$. La contención
opera en un solo sentido: un rol otorgado sobre un colegio alcanza sus sedes, pero ser
acudiente de una matrícula no confiere autoridad sobre la sede.

\paragraph{La función falla cerrada.}

Las políticas de seguridad a nivel de fila de los microservicios anteriores iniciaban su
condición base verificando si el contexto de usuario era nulo, de modo que toda conexión sin
contexto establecido veía la tabla completa. El control fallaba \emph{abierto}, y su modo de
fallo más probable, omitir el contexto, era precisamente el que concedía acceso total. La
función aquí especificada no tiene ninguna rama que conceda por ausencia de información.

En el sistema construido, identidad guarda las asignaciones de rol con su vigencia y los
pares de rol y permiso, y evalúa los permisos con ámbito (ADR 0013); las rutas protegidas
de inteligencia fallan cerradas (SG-06). La aplicación uniforme (AD-03 y SG-02) se delimitó como
endurecimiento productivo (Sección~\ref{sec:trabajo-futuro}): solo una ruta de producción consulta la función de permiso, y varias rutas que
devuelven datos personales no exigen sesión ni rol. Tampoco hay políticas de seguridad a
nivel de fila, de modo que la evaluación local del ADR 0016 sigue siendo diseño
(Sección~\ref{sec:defectos}).

\paragraph{Dos ejes independientes, no uno.}

La seguridad a nivel de campo usa el mismo mecanismo, con un componente de campo en el
permiso: leer un estudiante y leer su nombre son permisos distintos. El rediseño inicial
afirmó que esos permisos se derivaban \emph{automáticamente} del metadato de clasificación de
datos, y el análisis del registro de tratamiento mostró que no es así. La clasificación mide
sensibilidad y el permiso expresa quién puede ver qué: son dos ejes independientes. De la
clasificación solo se deriva el valor por omisión, de modo que todo campo clasificado como
dato de menor o como personal sensible nace denegado y todo campo público nace permitido.
Colapsar ambos ejes conduce a implementar la mitad del mecanismo creyendo que está completo.

% =====================================================================
\subsubsection{Protección de datos personales desde el diseño}


El sistema trata datos de menores de edad, datos personales de los acudientes, como el
contacto y la dirección de entrega, y datos de las instituciones educativas. La Ley 1581 de
2012 \citep{congreso_ley1581_2012} rige los dos primeros, y los terceros cuando identifican
a personas naturales: exige autorización previa y una finalidad legítima informada al
titular, y reconoce el derecho de supresión.

Se adoptó un esquema de cinco niveles de clasificación (público, interno, personal,
personal sensible y dato de menor), aplicado campo por campo en el registro de tratamiento
(Anexo~M), que cubre las cinco bases del sistema construido. Su primera versión, sobre los
microservicios anteriores, revisó 398 columnas de 50 tablas y encontró 126 campos con datos
personales, 20 de ellos clasificados como dato de menor.

\paragraph{Políticas que el esquema anterior no podía sostener.}

El ejercicio reveló un patrón recurrente: políticas correctamente enunciadas en la
documentación que el esquema hacía inejecutables. La Tabla~\ref{tab:politicas} recoge los
casos y su corrección.

% Tabla tab:politicas movida al Anexo A.3 (tablas del cuerpo)


El patrón tiene valor metodológico más allá de este proyecto: una política de tratamiento
de datos es verificable solo si el modelo lógico expone las columnas que la hacen
computable. Documentar la política sin instrumentar el esquema produce cumplimiento
aparente. El riesgo inverso también existe: el sistema construido tiene esas columnas, pero
la supresión propagada existe en el modelo y no como flujo de extremo a extremo
(Sección~\ref{sec:defectos}). El SRS v6.1 deja la supresión propagada, el panel de privacidad
y la retención automática fuera del alcance de este trabajo, y atiende los derechos del
titular mediante un caso de soporte.

\paragraph{Minimización por columna: la dirección de entrega.}

Los microservicios anteriores no tenían ninguna dirección de entrega, en un negocio cuya
propuesta de valor consiste en entregar el kit en el domicilio. Tampoco figuraba en el
esquema objetivo, de modo que la omisión se arrastraba del código al diseño: un defecto por
omisión sobrevive más que uno por comisión, porque nada lo señala.

El dato es sensible de forma poco habitual: más que información de contacto, es la
ubicación física del hogar de una familia con menores, cruzada con el colegio al que
asisten. El ADR 0015 distingue dos cosas que se parecen sin ser la misma: la libreta de
direcciones reutilizables del acudiente, en identidad, editable y sujeta a supresión, y la
copia congelada en el envío, en comercio, inmutable una vez confirmado el pedido. El
razonamiento es análogo al de congelar el precio en la línea del pedido: un hecho histórico
y un dato editable son entidades distintas aunque se escriban igual. Si el envío
referenciara la libreta, corregir un número de apartamento reescribiría el destino de una
entrega pasada.

La partición de columnas responde al principio de minimización. La ciudad y la localidad
quedan en claro porque de ellas depende la asignación de la zona de reparto, mientras que
la calle, el número y el teléfono se cifran en la aplicación: planificar una ruta requiere
conocer la zona, no el apartamento. En el sistema construido esa partición rige para la
libreta, que guarda cifrados la dirección y el contacto. La copia del envío conserva el
destinatario y la dirección en claro, como prueba ante un reclamo de la Ley 1480 de 2011. Desde el
ADR 0049 guarda además la latitud y la longitud del punto confirmado en el mapa, ambas nulables:
los pedidos anteriores y los que usan Google quedan sin coordenadas. El ADR 0044 había descartado
guardarlas por los términos de Google, mientras que la licencia ODbL de OpenStreetMap sí lo
permite. La supresión del titular borra las dos columnas y la firma del pedido las incluye, para
que un cambio posterior se note. Aun así, la minimización por columna no cubre todavía el dato
congelado.
